\documentclass[11pt, a4paper]{article}

\usepackage[utf8]{inputenc}
\usepackage[T1]{fontenc}
\usepackage[spanish]{babel}
\usepackage[margin=2.5cm]{geometry}
\usepackage{parskip}
\usepackage{booktabs}
\usepackage{titlesec}
\usepackage{hyperref}

\titlespacing*{\section}{0pt}{1em}{0.5em}
\titlespacing*{\subsection}{0pt}{0.8em}{0.3em}

\title{
    \textbf{Propuesta Técnica} \\[0.3cm]
    \large Digitalización e Indexación de Registros Históricos\\
    de Emigración a Cuba desde A Coruña
}
\author{
    Marcelo Ferreiro Sánchez\\
    \small \href{mailto:marcelo.fsanchez@udc.es}{marcelo.fsanchez@udc.es} \\
    \small +34 604 023 998
    \and
    José Romero Conde\\
    \small \href{mailto:j.rconde@udc.es}{j.rconde@udc.es} \\
    \small +34 658 755 850
}
\date{\today}

\begin{document}

\maketitle


% -----------------------------------------------
\section*{Pipeline Técnico}
% -----------------------------------------------

\textbf{1. Digitalización.} Los documentos serán escaneados a alta resolución,
generando archivos master que se conservarán localmente.

\textbf{2. Reconocimiento de texto manuscrito (HTR).} Se empleará un sistema de
HTR adaptado específicamente a la tipografía y características de los
manuscritos de época, garantizando una extracción fiable de los datos de cada
registro. 

\textbf{3. Estructuración y base de datos.} La información extraída se
almacenará en una base de datos relacional, vinculando cada campo con su
localización exacta en el documento original. Se incluirá una versión comprimida
del escaneo para consulta web.

\textbf{4. Plataforma web.} Se desarrollará una interfaz pública de consulta que permitirá 
buscar registros por nombre, fecha u otros campos, y visualizar el documento original asociado.


% -----------------------------------------------
\section*{Equipo}
% -----------------------------------------------

El equipo está formado por Marcelo Ferreiro y José Romero, ambos estudiantes en
último curso del Grado en Inteligencia Artificial por la Universidade da Coruña.


% -----------------------------------------------
\section*{Requisitos Hardware, Software y de Cloud}
% -----------------------------------------------

\subsection*{Hardware}
El proyecto requiere cierta potencia de cómputo debido a la muy posible
necesidad de aplicar un ajuste fino (\textit{fine tuning}) a los sistemas de HTR,
lo que requiere de una unidad de procesamiento gráfico con suficiente memoria
dedicada (más de 16GB), como una NVIDIA RTX 5080 o similar. 

Por otra parte, el otro requisito hardware presente, más crítico, es el
almacenamiento. Aún sin conocer con precisión la dimensión de los registros,
es necesario asegurar un guardado fiable, seguro y \textbf{redundante} de los
datos escaneados y extraídos. Estos dos requisitos requieren de la disposición
de un ordenador con una GPU de las características mencionadas, y capacidad
suficiente con redundancia local, probablemente con dos discos completamente
redundantes en espejo (RAID 1). Un equipo fiable que reúna tales características
puede rondar los 4.000€, incluyendo todos los periféricos necesarios para su uso.

Además, para el desarrollo del proyecto y conexión con la torre princial (el
equipo anteriormente mencionado) convendrán dos equipos portátiles. Por ejemplo
dos \textit{Macbook Air} básicos, rondando los 1100€ cada unidad aproximadamente.

En el año posterior al comienzo del proyecto podría evaluarse la necesidad de un
equipo de almacenamiento extra (un NAS) para disponer de redundancia en dos
equipos físicos separados. No sería necesario durante el primer año, ya que no se
dispondría aún de una cantidad de datos que lo haga necesario. Su coste es
considerablemente menor al del resto del hardware mencionado (menos de 800€).

Por último, quedando aún por decidir el método de digitalización, puede ser
necesario un escáner de formato A3 de alta resolución (más de 600 DPI). En tal
caso será necesario el uso de un escáner tipo \textit{flatbed}, de V o de tipo
cenital, que respete los lomos de los registros. Un equipo de tales
características, como el Scansnap SV600 ronda los 700€.

\subsection*{Software}
Se estiman los requisitos software como nulos, ya que se plantea el uso de
herramientas con licencias abiertas en todas las partes del proyecto.

\subsection*{Cloud}
Quedando por establecer el cómo y dónde se va a hospedar la plataforma, es necesario
estimar el coste de mantenimiento en la nube. Utilizando servicios como Azure, AWS
o Hetzner, el coste de mantener tanto la página web de consultas como la base de
datos puede estimarse en unos 100€ mensuales, dependiendo del tamaño del archivo,
que queda por determinar.

Este gasto en infraestructura en la Nube no sería necesario en caso de que sea
posible mantener la página y la base de datos en los servidores de la Xunta, en
tal caso su costo sería marginal. 

% -----------------------------------------------
\section*{Presupuesto}
% -----------------------------------------------

\begin{table}[h]
\centering
\begin{tabular}{@{}lr@{}}
\toprule
\textbf{Concepto} & \textbf{Importe anual} \\
\midrule
Honorarios (Marcelo Ferreiro)          & 10.000 € \\
Honorarios (José Romero)               & 10.000 € \\
Hardware (año 1: torre, escáner y equipos portátiles)      &  7.900 € \\
Hardware (año 2 estimado: NAS)         &    800 € \\
Software                               &      0 € \\
Servicios Cloud(en caso de ser necesario)                      &  1.200 € \\
\midrule
\textbf{Total año 1}                   & \textbf{29.100 €} \\
\textbf{Total año 2 y sucesivos}       & \textbf{22.000 €} \\
\bottomrule
\end{tabular}
\end{table}

\end{document}
